\section{Normalization chain for the square-root product}

This section records the constants in the normalization used in the
manuscript.  The claims below are identities after the external scalar
$C\in\C^\times$ in the protected curve-counting series has been fixed.
The Borcherds product fixes all other constants.

\subsection{Variables and Jacobi input}

Put
$$
Z=\begin{pmatrix}z_1&z_2\\ z_2&z_3\end{pmatrix},\qquad
q=e^{2\pi i z_1},\qquad r=e^{2\pi i z_2},\qquad
s=e^{2\pi i z_3}.
$$
Then
$$
\exp\bigl(\pi i(nz_1+lz_2+mz_3)\bigr)
=q^{n/2}r^{l/2}s^{m/2},
$$
while the Borcherds product factors use
$$
\exp\bigl(2\pi i(nz_1+lz_2+mz_3)\bigr)
=q^nr^ls^m.
$$

The Jacobi form used as Borcherds input is not the K3 elliptic genus
itself but half of it:
$$
Z_{\mathrm{K3}}(\tau,z)=2\phi_{0,1}(\tau,z),
\qquad
\phi_{0,1}(\tau,z)=\sum_{n\geq0,\;l\in\Z}f(n,l)q^nr^l.
$$
In the Eichler--Zagier normalization \cite{EZ:1},
$$
\phi_{0,1}(\tau,z)
=4\left(
\frac{\vartheta_2(\tau,z)^2}{\vartheta_2(\tau,0)^2}
+\frac{\vartheta_3(\tau,z)^2}{\vartheta_3(\tau,0)^2}
+\frac{\vartheta_4(\tau,z)^2}{\vartheta_4(\tau,0)^2}
\right),
$$
and the first terms are
$$
\phi_{0,1}
=(r^{-1}+10+r)
+q(10r^{-2}-64r^{-1}+108-64r+10r^2)+O(q^2).
$$
Thus
$$
f(0,-1)=1,\qquad f(0,0)=10,\qquad f(0,1)=1.
$$

\subsection{Primary-source variable map}

Oberdieck--Pixton \cite[Theorem 1]{OB} use
$$
P:=p_{\mathrm{OP}}=e^{iu}=e^{2\pi iz},\qquad
Q:=q_{\mathrm{OP}}=e^{2\pi i\tau},\qquad
T:=\widetilde q_{\mathrm{OP}}=e^{2\pi i\widetilde\tau}.
$$
Their primitive Gromov--Witten partition function is
$$
\mathcal Z_{\mathrm{OP}}(u,Q,T)
=\sum_{g,h,d}
N_{g,h,d}u^{2g-2}Q^{h-1}T^{d-1}.
$$
Oberdieck--Pandharipande use the same Siegel variables
\cite[Section 2]{OPand}:
$$
\Omega=\begin{pmatrix}\tau&z\\ z&\widetilde\tau\end{pmatrix},
\qquad
P=e^{2\pi iz},\quad Q=e^{2\pi i\tau},\quad T=e^{2\pi i\widetilde\tau}.
$$
The Borcherds variables of this manuscript are
$$
q=Q,\qquad r=P,\qquad s=T.
$$
The Hilbert-scheme Donaldson--Thomas variables in the displayed series
$$
Z^X_{\mathrm{DT}}(q,t,p)
=\sum_{h,d,n}\mathbf{DT}^X_{n,(\beta_h,d)}
q^{d-1}t^{h-1}(-p)^n
$$
are therefore
$$
q_{\mathrm{DT}}=T=s,\qquad
t_{\mathrm{DT}}=Q=q,\qquad
p_{\mathrm{DT}}=P=r.
$$
The last equality uses Bryan's convention \cite[Definition 2.1]{BRYAN}
$$
\mathbf{DT}(X)=
\sum_{h,d,n}\mathbf{DT}_{h,d,n}(X)\,
\widetilde q^{\,h-1}q^{d-1}(-p)^n
$$
and Bryan's warning that his $q,\widetilde q$ are opposite to the
Oberdieck--Pandharipande convention.  The stable-pairs sign is
compatible with Oberdieck--Pandharipande's convention
\cite[Section 2]{OPand}
$$
y=-P.
$$

\begin{lemma}[Weight fixed by the zero coefficient]
The Borcherds lift of the input $\phi_{0,1}$ has weight
$$
\frac{f(0,0)}{2}=5.
$$
If one uses $Z_{\mathrm{K3}}=2\phi_{0,1}$ as input instead, all product
exponents and the weight double; the resulting scalar genus-two form is
the square $\Delta_5^2=\Delta_{10}$.
\end{lemma}

\begin{proof}
Borcherds' weight formula for the multiplicative lift of a weak Jacobi
form of weight zero and index one gives weight one half of the
coefficient of $q^0r^0$ \cite{B}, \cite{GN}.  The displayed expansion gives
$f(0,0)=10$, hence the weight is $5$.  Replacing the input by
$2\phi_{0,1}$ replaces each coefficient $f(n,l)$ by $2f(n,l)$, so the
weight becomes $10$ and the multiplicative product is squared.
\end{proof}

\subsection{The coefficient 64}

The Igusa normalization in this manuscript is fixed by the product of the
ten even genus-two theta constants \cite{FREITAG:1}, \cite{VDGEER:1},
\cite{MAASS:1}
$$
\Delta_5(Z)=
\prod_{\substack{a,b\in(\Z/2\Z)^2\\ a\cdot b\equiv0\!\!\pmod2}}
\nu_{a,b}(Z),
$$
where
$$
\nu_{a,b}(Z)=
\sum_{l\in\Z^2}
\exp\bigl(\pi i(Z[l+a/2]+b\cdot l)\bigr).
$$
Equivalently, for $u_i=l_i+a_i/2$,
$$
\nu_{a,b}(Z)=
\sum_{l\in\Z^2}
(-1)^{b_1l_1+b_2l_2}
q^{u_1^2/2}r^{u_1u_2}s^{u_2^2/2}.
$$

The leading terms by the first characteristic $a$ are:
$$
\begin{array}{c|c|c}
a & \hbox{allowed } b & \hbox{lowest product in this } a\hbox{-block}\\
\hline
(0,0) & (0,0),(1,0),(0,1),(1,1) & 1\\
(1,0) & (0,0),(0,1) & 4q^{1/4}\\
(0,1) & (0,0),(1,0) & 4s^{1/4}\\
(1,1) & (0,0),(1,1) &
4q^{1/4}s^{1/4}(r^{1/2}-r^{-1/2})
\end{array}
$$
The last row is the only row with a nontrivial $r$-sign:
$$
\nu_{(1,1),(0,0)}
=2q^{1/8}s^{1/8}(r^{1/4}+r^{-1/4})+\cdots,
$$
$$
\nu_{(1,1),(1,1)}
=2q^{1/8}s^{1/8}(r^{1/4}-r^{-1/4})+\cdots.
$$
Multiplication gives
$$
\Delta_5(Z)
=64q^{1/2}s^{1/2}(r^{1/2}-r^{-1/2})+\cdots
$$
and hence, in the Fourier convention
$$
\Delta_5(Z)
=\sum c_\Delta(n,l,m)\exp\bigl(\pi i(nz_1+lz_2+mz_3)\bigr),
$$
one has
$$
c_\Delta(1,1,1)=64,\qquad c_\Delta(1,-1,1)=-64.
$$
Thus the sign of $\Delta_5$ is fixed by the condition
$c_\Delta(1,1,1)=64$.

\subsection{Borcherds product, vacuum factor, and 1/64}

Let
$$
\Gamma_{\mathrm{eff}}
=\{(n,l,m)\mid m>0,\ n\geq0,\ l\in\Z\}
\cup\{(n,l,0)\mid n>0,\ l\in\Z\}
\cup\{(0,l,0)\mid l<0\}.
$$
The Gritsenko--Nikulin Borcherds product attached to $\phi_{0,1}$ is
$$
D_5(Z):=\frac{1}{64}\Delta_5(Z)
=q^{1/2}r^{1/2}s^{1/2}
\prod_{(n,l,m)\in\Gamma_{\mathrm{eff}}}
(1-q^nr^ls^m)^{f(nm,l)}.
$$
The word ``monic'' refers to the vacuum monomial
$q^{1/2}r^{1/2}s^{1/2}$ before expansion of the real-root factor
$(1-r^{-1})^{f(0,-1)}=(1-r^{-1})$.  Expanding that factor gives
$$
D_5(Z)
=q^{1/2}s^{1/2}(r^{1/2}-r^{-1/2})+\cdots,
$$
which is exactly the theta-constant computation above divided by $64$.

The determinant used in the body of the paper is the non-monic Igusa
normalization
$$
\mathcal D_X(Z):=\Delta_5(Z)
=64D_5(Z).
$$
Equivalently, with
$$
v_0(Z):=64q^{1/2}r^{1/2}s^{1/2},
$$
one has
$$
\mathcal D_X(Z)
=v_0(Z)
\prod_{(n,l,m)\in\Gamma_{\mathrm{eff}}}
(1-q^nr^ls^m)^{f(nm,l)}.
$$
For the one-particle superspace
$$
\mathcal V=\bigoplus_{\gamma\in\Gamma_{\mathrm{eff}}}\mathcal V_\gamma,
\qquad
\sdim\mathcal V_{(n,l,m)}=f(nm,l),
$$
the Fock character identity gives
$$
\sch\mathcal F(\mathcal V)
=\prod_{(n,l,m)\in\Gamma_{\mathrm{eff}}}
(1-q^nr^ls^m)^{-f(nm,l)}.
$$
Therefore
$$
\mathcal Z_{\mathrm{BPS}}^{\mathrm{vir}}(Z)
:=\mathcal D_X(Z)^{-1}
=\frac1{64}q^{-1/2}r^{-1/2}s^{-1/2}\sch\mathcal F(\mathcal V).
$$
This is the source of the factor $1/64$: it is the inverse of the Igusa
leading coefficient, not a Fock-space multiplicity.

\subsection{The doubling substitution}

Let
$$
z=z_1f_{-2}+z_2f_3+z_3f_2,
\qquad
\alpha(n,l,m)=2nf_2-lf_3+2mf_{-2}.
$$
With the lattice pairing used in the manuscript,
$$
q^nr^ls^m=\exp\bigl(-\pi i(\alpha(n,l,m),z)\bigr).
$$
The Weyl--Kac--Borcherds denominator is written with
$\exp(-2\pi i(\,\cdot\,,z))$.  Consequently the manuscript passes from
the Borcherds product convention to the denominator convention by the
substitution
$$
Z\longmapsto2Z:
\qquad
\exp\bigl(-\pi i(\lambda,z)\bigr)
\longmapsto
\exp\bigl(-2\pi i(\lambda,z)\bigr).
$$
With $\rho=f_2-\frac12f_3+f_{-2}$,
$$
D_5(2Z)
=\frac1{64}\Delta_5(2Z)
=e^{-2\pi i(\rho,z)}
\prod_{\alpha\in\mathcal R_+}
\left(1-e^{-2\pi i(\alpha,z)}\right)^{\smult(\alpha)}.
$$
Thus
$$
\operatorname{den}(\mathfrak g_{\Delta_5})
=\frac1{64}\Delta_5(2Z).
$$

\subsection{The scalar square and the constant C}

The character of $\Delta_5$ squares to the trivial character, and the
scalar Igusa cusp form is
$$
\Delta_{10}(Z)=\Delta_5(Z)^2.
$$
In the present normalization this is
$$
\Delta_{10}(Z)
=4096\,qrs
\prod_{(n,l,m)\in\Gamma_{\mathrm{eff}}}
(1-q^nr^ls^m)^{2f(nm,l)}.
$$
Equivalently,
$$
\frac1{4096}\Delta_{10}(Z)=D_5(Z)^2.
$$

Oberdieck--Pixton define \cite[(3)]{OB}
$$
\chi_{10}^{\mathrm{OP}}(P,Q,T)
=PQT
\prod_{k,h,d}
(1-P^kQ^hT^d)^{c_{\mathrm{K3}}(4hd-k^2)},
$$
where the product runs over
$$
k\in\Z,\qquad h,d\geq0,\qquad
(h>0\ \hbox{or}\ d>0)\quad\hbox{or}\quad(h=d=0,\ k<0).
$$
The coefficient function is the K3 elliptic genus:
$$
Z_{\mathrm{K3}}(P,Q)
=\sum_{h\geq0,\;k\in\Z}
c_{\mathrm{K3}}(4h-k^2)P^kQ^h
=2\phi_{0,1}(Q,P).
$$
Hence
$$
c_{\mathrm{K3}}(4hd-k^2)=2f(hd,k),
$$
and the OP product is the square of the monic Borcherds product:
$$
\chi_{10}^{\mathrm{OP}}(P,Q,T)=D_5(Z)^2
=4096^{-1}\Delta_5(Z)^2.
$$
Oberdieck--Pixton prove \cite[Theorem 1]{OB}
$$
\mathcal Z_{\mathrm{OP}}(u,Q,T)
=-\left(\chi_{10}^{\mathrm{OP}}(P,Q,T)\right)^{-1},
\qquad P=e^{iu}.
$$
Therefore, in the theta-constant normalization of this manuscript,
$$
\mathcal Z_{\mathrm{OP}}(u,Q,T)
=-4096\,\Delta_5(Z)^{-2}.
$$
The sign is the sign in the primary theorem; the factor $4096$ is
$64^2$ and comes only from replacing the monic product $D_5$ by
$\Delta_5$.

For a general scalar-square branch the partition function is normalized
only up to a nonzero scalar $C$:
$$
Z^X_{\square}(q,r,s)=\frac{C}{\Delta_5(Z)^2}
=\frac{C}{\Delta_{10}(Z)}.
$$
For the OP branch,
$$
C_{\mathrm{OP}}=-4096.
$$
For the Hilbert-scheme DT branch displayed above, no primary source
used here fixes the equality with the OP branch in all degrees.  The
honest statement is the assumption
$$
Z^X_{\mathrm{DT}}(q_{\mathrm{DT}},t_{\mathrm{DT}},p_{\mathrm{DT}})
=C_{\mathrm{DT}}\Delta_5(Z)^{-2},
\qquad C_{\mathrm{DT}}\in\C^\times.
$$
Hence
$$
\frac{C}{Z^X_{\square}}=\Delta_5(Z)^2=\Delta_{10}(Z),
\qquad
\sqrt{\frac{C}{Z^X_{\square}}}=\Delta_5(Z),
$$
where the square-root branch is the branch with leading vacuum factor
$$
64q^{1/2}r^{1/2}s^{1/2}.
$$
Consequently
$$
\mathcal Z_{\mathrm{BPS}}^{\mathrm{vir}}=\Delta_5^{-1},
\qquad
\left(\mathcal Z_{\mathrm{BPS}}^{\mathrm{vir}}\right)^2
=\frac{Z^X_{\square}}{C}.
$$
The scalar $C$ belongs to the enumerative normalization of
$Z^X_{\square}$; it is not produced by the Borcherds product.

\subsection{Audit outcome}

The normalizations fixed by primary sources are:
$$
\begin{array}{c|c|c}
\hbox{normalization} & \hbox{value} & \hbox{source}\\
\hline
\hbox{OP variables} &
(P,Q,T)=(e^{iu},q_{\mathrm{OP}},\widetilde q_{\mathrm{OP}})
& \cite[Theorem 1]{OB}\\
\hbox{OP partition function} &
\mathcal Z_{\mathrm{OP}}=-1/\chi_{10}^{\mathrm{OP}}
& \cite[Theorem 1]{OB}\\
\hbox{OP product} &
\chi_{10}^{\mathrm{OP}}=D_5^2
& \cite[(3)]{OB}\\
\hbox{this manuscript} &
D_5=64^{-1}\Delta_5
& \cite{GN},\cite{B}\\
\hbox{constant conversion} &
\chi_{10}^{\mathrm{OP}}=4096^{-1}\Delta_5^2
& c_\Delta(1,1,1)=64\\
\hbox{OP scalar} &
C_{\mathrm{OP}}=-4096
& -1/\chi_{10}^{\mathrm{OP}}\\
\hbox{DT variables} &
(q_{\mathrm{DT}},t_{\mathrm{DT}},p_{\mathrm{DT}})=(T,Q,P)
& \cite[Definition 2.1]{BRYAN}
\end{array}
$$
The remaining ambiguity is not a sign or a factor of $64$.  It is the
unproved identification
$$
Z^X_{\mathrm{DT}}
\stackrel{?}{=}
\mathcal Z_{\mathrm{OP}}
$$
in the Behrend-weighted Hilbert-scheme normalization.  Until that
comparison is supplied, the manuscript must keep $C_{\mathrm{DT}}$
separate from $C_{\mathrm{OP}}=-4096$.

\subsection{Check table}

$$
\begin{array}{c|c|c}
\hbox{object} & \hbox{normalization} & \hbox{consequence}\\
\hline
q,r,s
& e^{2\pi iz_1},e^{2\pi iz_2},e^{2\pi iz_3}
& \exp(\pi i(nz_1+lz_2+mz_3))=q^{n/2}r^{l/2}s^{m/2}\\
Z_{\mathrm{K3}}
& Z_{\mathrm{K3}}=2\phi_{0,1}
& \hbox{Borcherds input is }\phi_{0,1}\\
\phi_{0,1}
& r^{-1}+10+r+O(q)
& f(0,0)=10\\
\hbox{weight}
& f(0,0)/2
& 5\\
\Delta_5
& \prod_{\mathrm{even}}\nu_{a,b}
& c_\Delta(1,1,1)=64\\
D_5
& 64^{-1}\Delta_5
& q^{1/2}r^{1/2}s^{1/2}\prod(1-q^nr^ls^m)^{f(nm,l)}\\
\mathcal D_X
& \Delta_5
& v_0=64q^{1/2}r^{1/2}s^{1/2}\\
\mathcal Z_{\mathrm{BPS}}^{\mathrm{vir}}
& \mathcal D_X^{-1}
& 64^{-1}q^{-1/2}r^{-1/2}s^{-1/2}\sch\mathcal F(\mathcal V)\\
\operatorname{den}(\mathfrak g_{\Delta_5})
& D_5(2Z)
& 64^{-1}\Delta_5(2Z)\\
\Delta_{10}
& \Delta_5^2
& 4096\,D_5^2\\
Z^X_{\mathrm{OP}}
& -4096\,\Delta_5^{-2}
& \hbox{Oberdieck--Pixton primitive theorem}\\
Z^X_{\square}
& C/\Delta_5^2
& (\mathcal Z_{\mathrm{BPS}}^{\mathrm{vir}})^2=Z^X_{\square}/C
\end{array}
$$
